\scanpage{254}
constant. En effet (\unsure{cf} ?) $H^i(X_n,\mathbb Z)=\unsure{H^i}(X_0,\mathbb Z)$ est un groupe de type fini, et son quotient par $\operatorname{Ker}\alpha^i$ est sans torsion [car isom à un sous-groupe de $H^i(X,\mathcal O_{X_n})$, qui est un vectoriel sur $\mathbb C$], \deleted{\unsure{d'où}} il suffit donc de noter que pour $n$ grand, le rang du dit quotient est constant.

On obtient ainsi, compte tenu de 1.3 (ii) \inserted{et du fait qu'une ext. d'un syst. M.L. par un syst. M.L. est M.L.}

\uline{\deleted{Corollaire}\inserted{Proposition}} 1.5 Sous les conditions de 1.3,

(i) le syst. projectif $(H^i(X_n,\mathcal O_{X_n}^{\unsure{\wedge|*}}))_n$ satisfait ML

(ii) On a une suite exacte
\diagram{\begin{tikzcd}[column sep=small]
H^i(X_0,\mathbb Z)\arrow[r,"\alpha_\infty^{(i)}"] & {R^if_*(\mathcal O_X)_y}^{\wedge}\arrow[r,"\beta_\infty^{(i)}"] & \varprojlim H^i(X_n,\mathcal O_{X_n}^{\unsure{\wedge|*}})\arrow[r,"\delta_\infty^{(i)}"] & H^{i+1}(X_0,\mathbb Z)\arrow[r,"\alpha_\infty^{(i+1)}"] & {R^{i+1}f_*(\mathcal O_X)_y}^{\wedge}
\end{tikzcd}}
\margin{\deleted{\unsure{*}}}

Notons qu'on a une suite exacte déduite de la suite exacte
\diagram{\begin{tikzcd}
0\arrow[r,dash] & \mathbb Z_X\arrow[r,dash] & \mathcal O_X\arrow[r] & \mathcal O_X^{\unsure{\wedge|*}}\arrow[r] & 0
\end{tikzcd}}
par application des $(R^if_*)_y$ et compte tenu de 1.2, les \inserted{\unsure{fibres en $y$ des}} termes $R^if(\unsure{\mathbb Z|\mathcal O}_X)_{\deleted{\unsure{y}}}$ \deleted{se} s'identifient à $H^i(X_0,\mathbb Z)$. On obtient ainsi une
\mspage{253}
